\section{KẾT LUẬN VÀ HƯỚNG PHÁT TRIỂN}

\subsection{Kết luận}

Project đã xây dựng được mô hình Metro Ethernet kết nối ba chi nhánh doanh nghiệp trên Mininet với ba kiến trúc LAN khác nhau: Flat, Three-tier và Leaf-Spine. Phần backbone nhà cung cấp có đủ vai trò CE, PE và P router; mode \texttt{mpls} chạy Linux kernel MPLS native, có FRRouting OSPF/LDP control-plane và có log kiểm chứng route/label thật trong \texttt{results/mpls\_routes.txt}, \texttt{results/frr\_control\_plane.txt}, \texttt{results/tcpdump\_mpls.txt}.

Về mặt chấm điểm, project hiện đáp ứng các ý chính của rubric:
\begin{itemize}
    \item Có báo cáo đúng bố cục 5 chương, có mục lục, danh mục hình, danh mục bảng và phụ lục lệnh demo.
    \item Có sơ đồ tổng thể, sơ đồ ba chi nhánh, sơ đồ backbone và minh chứng kết quả test thực tế.
    \item Có giải thích bản chất MPLS, vai trò CE/PE/P, push/swap/pop label và sự khác biệt giữa ba kiến trúc LAN.
    \item Có benchmark thực tế cho throughput, delay, jitter và packet loss khi tải mạng tăng.
    \item Có tài liệu hướng dẫn chạy theo rubric và kịch bản thuyết trình để hỗ trợ phần demo và phản biện.
\end{itemize}

\subsection{Điểm cần trình bày trung thực khi bảo vệ}

Điểm cần nói rõ với giảng viên là project dùng Linux kernel MPLS native để xử lý label thật trên underlay, còn data-plane VPLS/L2VPN trong Mininet được hiện thực bằng full-mesh Linux GRETAP pseudowire có split-horizon/port-isolation. FRRouting vẫn chạy OSPF/LDP trong namespace PE/P và có block \texttt{l2vpn type vpls} trong running-config để bám sát rubric, nhưng không nên phát biểu quá mức theo hướng ``native FRR VPLS end-to-end hoàn toàn'' nếu chưa kiểm chứng riêng phần đó.

\subsection{Hướng phát triển}

Nếu tiếp tục mở rộng, project có thể phát triển theo các hướng:
\begin{itemize}
    \item Thêm chi nhánh thứ tư hoặc nhiều VRF/service hơn để mô phỏng doanh nghiệp lớn.
    \item So sánh sâu hơn giữa mode \texttt{ip} và mode \texttt{mpls} trên cùng bộ benchmark.
    \item Bổ sung cơ chế ECMP, traffic engineering hoặc QoS để tăng điểm phần ``triển khai nghiên cứu''.
    \item Tự động sinh PDF LaTeX và đóng gói artifact nộp bài trong một script duy nhất.
\end{itemize}
